\documentclass{article}

\usepackage[preprint]{neurips_2024}

\usepackage[utf8]{inputenc}
\usepackage[T1]{fontenc}
\usepackage{hyperref}
\usepackage{url}
\usepackage{booktabs}
\usepackage{amsfonts}
\usepackage{amssymb}
\usepackage{nicefrac}
\usepackage{microtype}
\usepackage{xcolor}
\usepackage{graphicx}
\usepackage{multirow}

\title{Are Polymarket Bettors Smarter than Economists?\\
       Head-to-Head on Inflation Forecasts}

\author{%
  Aly Sultan \\
  Independent Researcher \\
  \texttt{aly.sultan@protonmail.com} \\
  \And
  Mei-Ling Wong \\
  Department of Economics \\
  University of Auckland \\
  \texttt{m.wong@auckland.ac.nz} \\
  \And
  Daniel Brüggemann \\
  Institute for Applied Microeconomics \\
  Goethe-Universit\"at Frankfurt \\
  \texttt{brueggemann@uni-frankfurt.de} \\
}

\begin{document}

\maketitle

\begin{abstract}
Polymarket attracted roughly \$18M of trading volume on the 29 US CPI
prediction markets that resolved between January 2025 and April 2026
in our sample.  Are these markets actually useful forecasts of
inflation, or just noise?  We compare the bracket-implied point
forecast from those 29 events, covering monthly (MoM) and annual
(YoY) US CPI prints, to the Federal Reserve Bank of Cleveland's daily
inflation nowcasting model.  In the
unadjusted head-to-head, the two forecasters are statistically
indistinguishable: both have a mean absolute error of 0.124 percentage
points (paired $t = 0.02$, $p = 0.99$).  Splitting by horizon, Polymarket
edges the model on monthly markets (MAE 0.080 vs 0.123\,pp) while the
model wins on annual markets (MAE 0.124 vs 0.160\,pp).  The headline tie
is fragile: a single event whose original bracket grid did not span the
realised value (March 2026 YoY, off by $-0.93$\,pp) accounts for almost
all of Polymarket's RMSE excess; on a sample restricted to the
operator-curated bracket grids that actually contained the realised
print, Polymarket has lower MAE than the Cleveland Fed model
($0.078$ vs $0.117$\,pp, $p = 0.10$ on a paired $t$-test).  We
conclude that liquid Polymarket inflation markets are roughly competitive
with a professional nowcasting system, with the caveat that small
samples and market-design choices (bracket coverage, addendum markets)
materially affect the answer.  All numbers are reproducible from the
attached scripts; the largest practical bottleneck was Polymarket's
14-day CLOB price-history window, which forces a short look-back.
\end{abstract}

\section{Introduction}

Prediction markets have been studied for decades as information
aggregators, with classic results showing they often match or beat
expert panels on political and sporting outcomes
\citep{wolfers2004prediction,arrow2008promise,berg2008election}.
Polymarket, a USDC-settled exchange running on the Polygon
blockchain, brought prediction markets back into the mainstream during
the 2024 US presidential election: by election day it had cleared an
estimated \$2.4B of volume on Trump-vs-Harris contracts alone
\citep{tsang2026polymarket}.

A natural question is whether Polymarket's apparent edge on political
horse-races extends to macroeconomic forecasts -- specifically to
Consumer Price Index (CPI) prints, which are difficult to forecast
\citep{stockwatson2007inflation,faust2013forecasting} and historically
the domain of professional surveys and structural models
\citep{ang2007macro,knotek2017nowcasting}.  Anecdotally, Polymarket
``Inflation -- Monthly'' and ``Inflation -- Annual'' markets routinely
attract \$0.5M -- \$3M in volume per print, suggesting both that
informed traders find these markets worth working and that the operator
finds them worth listing.  Whether the resulting prices are actually
informative has not, to our knowledge, been measured in a published
academic comparison.

We perform such a comparison, with two questions in mind:

\begin{enumerate}
\item How accurate are Polymarket-implied CPI forecasts at the
      typical T-1-day horizon?
\item How do those forecasts compare to a transparent professional
      benchmark -- the Federal Reserve Bank of Cleveland's daily
      inflation nowcasting model?
\end{enumerate}

Our contribution is a small but real-data, fully reproducible
head-to-head on 29 paired events covering target months between
January 2025 and April 2026.  The headline result is a statistical tie
which dissolves into a Polymarket-favouring result once a single
truncated-bracket event is handled with the operator's own
``Higher Brackets'' addendum.  We also document the market-design
artefact -- when the original bracket grid does not cover the realised
print, the bracket-implied mean is biased toward the bracket boundary
-- and quantify its impact.

\textbf{Scope and honesty.}  Polymarket's public CLOB API limits
historical price queries to roughly 14-day windows at 1-hour fidelity,
which constrains us to short-horizon ``T - 1 day'' forecasts.  We have
no access to a survey-based ``true'' economist consensus (the
Philadelphia Fed Survey of Professional Forecasters is quarterly and
does not nowcast individual monthly prints), so the ``economist'' in
our title is the Cleveland Fed nowcast -- a model, run by economists,
publicly documented in \citet{knotek2017nowcasting}.  Sample sizes are
small ($N \leq 29$) and effect sizes are modest; we are explicit about
this throughout.

\section{Related Work}

\paragraph{Prediction markets as forecasters.}
\citet{wolfers2004prediction} survey the theory and empirical record of
prediction markets, showing that market-generated forecasts are usually
at least as accurate as comparable benchmarks.  The market scoring rule
machinery underpinning many modern prediction markets, including the
LMSR-style automated market makers, was developed by
\citet{hanson2003combinatorial}.  \citet{snowberg2013prediction} review
the macroeconomic-forecasting record specifically, finding that
prediction markets ``quickly incorporate new information, are largely
efficient, and \dots\ exhibit lower statistical errors than
professional forecasters and polls''.  Empirical assessments are mixed:
\citet{atanasov2017distilling} find that prediction markets beat the
simple mean of polls but lose to algorithmically aggregated polls with
recalibration; \citet{page2013calibration} document a horizon-dependent
favourite/longshot bias.  Theoretical interpretation of market prices
as probabilities is non-trivial -- \citet{manski2006interpreting} shows
the equilibrium price under heterogeneous beliefs is the
\emph{midpoint} of an interval containing the mean belief, not the
mean itself, and \citet{tetlock2008liquidity} documents that liquidity
does not monotonically improve efficiency.

\paragraph{Polymarket.}
\citet{tsang2026polymarket} use Polymarket's on-chain transaction
record from the 2024 US presidential election to separate genuine
volume from share creation, and find that Polymarket's effective
trading volume was less than half of the headline figure but that
market efficiency (measured by Kyle's lambda and arbitrage delay)
improved substantially over the campaign.

\paragraph{Inflation forecasting.}
\citet{stockwatson2007inflation} document a structural break in the
forecastability of US inflation around the mid-1980s and propose a
trend-cycle model with stochastic volatility.
\citet{ang2007macro} find that survey-based measures outperform
ARIMA, term-structure, and Phillips-curve models for US inflation
forecasts.  \citet{faust2013forecasting} report that simple
glide-path methods are hard to beat in the long run, while at short
horizons sophisticated nowcasting wins.  \citet{knotek2017nowcasting}
describe the Cleveland Fed's nowcasting model that we use as our
economist benchmark; that model has been continuously published on the
Federal Reserve Bank of Cleveland website each business day.
\citet{croushore2001realtime} and \citet{faust2005news} emphasise the
importance of using real-time, point-in-time data for honest
out-of-sample comparisons, which we do here.

\section{Data}

\paragraph{Polymarket events.}  We query Polymarket's public Gamma API
(\url{https://gamma-api.polymarket.com/events}) with the tag slugs
\texttt{cpi} and \texttt{inflation} (deduplicated by event id),
filtering to US monthly (MoM) and annual (YoY) CPI markets.  This
yields 43 events with end dates between December 2023 and May 2026.
Each event is a collection of mutually-exclusive
\emph{bracket markets}.  For a typical event (May 2025 YoY) the
bracket markets are ``Will annual inflation increase by 2.2\% or less
in May?'', ``\ldots by 2.3\%?'', \ldots,
``\ldots by 2.6\% or more in May?'' -- five brackets including
$\le$ and $\ge$ tail brackets at the boundaries.  Some events use
only point (\textsf{eq}) brackets with no tail buckets, which we
discuss further in Section~5.3.  We parse the bracket
specification from each market's question with a regex (see
\texttt{scripts/fetch\_data.py}).  For each bracket market we then fetch
the historical YES-token price series from Polymarket's CLOB API
(\url{https://clob.polymarket.com/prices-history}) at 1-hour fidelity.
The CLOB caps queries at a roughly 15-day span, so we anchor each query
to the 14 days immediately preceding the event end date -- which is
also the BLS release date for the target CPI print.

\paragraph{Cleveland Fed nowcasts.}  The Federal Reserve Bank of
Cleveland publishes daily MoM and YoY inflation nowcasts at
\url{https://www.clevelandfed.org/indicators-and-data/inflation-nowcasting},
along with the eventual BLS-released realised value once available.
The underlying time-series live in two JSON ``webchart'' files
(\texttt{nowcast\_year.json} and \texttt{nowcast\_month.json}); we
mirror these once per run.  Each ``chart'' in those files corresponds
to a target month and contains, as parallel time series, the model's
day-by-day nowcast (CPI, Core CPI, PCE, Core PCE) and the eventual
``Actual CPI Inflation'' value (sourced from the BLS press release).
We use the headline CPI series throughout, MoM for the
\emph{Inflation -- Monthly} markets and YoY for the
\emph{Inflation -- Annual} markets.

\paragraph{Sample.}  Joining the Polymarket events to Cleveland Fed
target-month charts leaves us with 38 events that have at least one
forecast available, of which 29 have both a Polymarket-implied point
forecast and a Cleveland Fed nowcast and a realised value -- our
paired-comparison sample.  Five events drop out: three early 2024
binary ``inflation > $X$\%?'' markets that pre-date the
bracketed-market era and whose question format our target-month parser
cannot handle, and two October 2025 events whose realised values have
not yet appeared in the Cleveland Fed chart we mirrored.\footnote{The
BLS October 2025 CPI release was delayed because of a US Federal
shutdown; this is also the subject of one of the open Polymarket
events (id 87073, ``BLS delays another CPI release before 2027?''),
which we do not analyse.}  Total Polymarket dollar volume across
the paired events is \$18.3M, with a median per-event volume of
\$591k.

\section{Methodology}

\subsection{Polymarket-implied point forecast}

For a given event $e$ with bracket markets $\{m_1, \ldots, m_K\}$ and
last-known YES probabilities $\{p_1, \ldots, p_K\}$ at the cutoff
timestamp, we recover the implied probability distribution over
bracket midpoints.  Let $b_k = (\text{kind}_k, v_k)$ denote the
bracket specification: $(\textsf{le}, v)$ for ``$\le v$''
(left-truncated lower-tail bracket), $(\textsf{eq}, v)$ for the
0.1\,pp-wide point bracket centred on $v$, and $(\textsf{ge}, v)$ for
``$\ge v$'' (right-truncated upper-tail bracket).  We assign midpoints
$m_k^{\star}$ as

\begin{equation*}
m_k^{\star} = \begin{cases}
v_k - 0.05 & \text{if kind} = \textsf{le} \\
v_k        & \text{if kind} = \textsf{eq} \\
v_k + 0.05 & \text{if kind} = \textsf{ge}
\end{cases}
\end{equation*}

\noindent
and renormalise the prices so $\sum_k p_k = 1$ before taking the
expectation:

\begin{equation*}
\hat{\pi}_{\text{poly}}(e) = \sum_{k=1}^{K}
   \frac{p_k}{\sum_{j} p_j} \cdot m_k^{\star}.
\end{equation*}

\noindent
The renormalisation absorbs the typical few-percent quote-spread sum
deviation from unity.  Following \citet{manski2006interpreting} we
emphasise that $\hat{\pi}_{\text{poly}}$ is the price-implied
expectation, not the expected belief of the marginal trader.

\paragraph{T-1 day forecast cutoff.}  CPI prints are released by the
BLS at 8:30 AM ET on the published release date.  Polymarket markets
resolve over the subsequent few hours.  We define the ``T-1 day''
forecast as the last hourly price recorded $\le 24$ hours before each
event's \texttt{endDate} field.

\subsection{Cleveland Fed nowcast at matched cutoff}

For each target month $(y, m)$ we look up the corresponding Cleveland
Fed chart (subcaption ``$y$-$m$''), iterate the daily ``CPI Inflation''
series in order, and return the last value with date $\le$ the
Polymarket T-1 day cutoff (using MM/DD string ordering, which works
because the Cleveland Fed chart always covers a single contiguous
window starting on the first business day of the target month).

\subsection{Metrics}

For each event, the forecast error is the signed difference
\[
\varepsilon^{(s)}_e = \hat{\pi}_s(e) - \pi^{\text{actual}}_e
\quad \text{for } s \in \{\text{poly}, \text{cf}\}.
\]
\noindent
We report mean absolute error
$\mathrm{MAE} = \tfrac{1}{N}\sum |\varepsilon|$,
root-mean-squared error $\mathrm{RMSE}$, signed bias
$\overline{\varepsilon}$, and Pearson and Spearman correlations
between forecast and realised values.  For Polymarket's full
probability distribution over brackets, we also compute the multi-class
Brier score \citep{brier1950verification}:
\[
B(e) = \sum_k (p_k - y_k)^2, \qquad
y_k = \mathbf{1}[\text{realised} \in \text{bracket } k].
\]
\noindent
$B$ ranges from 0 (perfect) to 2 (all mass on the wrong bracket); a
uniform predictor over 5 brackets has expected $B = 0.80$.

We test the headline hypothesis ``Polymarket and Cleveland Fed have
equal MAE'' with a paired $t$-test on $\{|\varepsilon^{(\text{poly})}_e|
- |\varepsilon^{(\text{cf})}_e|\}$.

\subsection{Robustness}

\textbf{Drop-1.}  We re-compute headline metrics after dropping the
single event with the largest absolute difference in $|\varepsilon|$
between forecasters.

\textbf{Bracket-addendum dedupe.}  Polymarket occasionally re-lists a
target with a wider bracket grid when the original grid looks
underspecified -- ``March Inflation US -- Annual (Higher Brackets)''
was added when the original grid's top bracket of $2.7\%$ was below
the realised print of $3.26\%$.  For each (target month, kind), we
keep the addendum event if one exists, otherwise the single original
event.  We report the dedupe-sample metrics as a robustness check;
this also helps mitigate the truncation bias documented below.

\section{Experiments and Results}

\subsection{Headline}

Table~\ref{tab:headline} summarises the head-to-head on all 29 paired
events.  Both forecasters have a mean absolute error of
0.124\,pp.  The paired $t$-test on absolute errors gives
$t = 0.02, p = 0.99$ -- a stark null.  Both forecasts are strongly
correlated with the realised CPI value: Pearson $r = 0.984$ for
Polymarket and $r = 0.992$ for the Cleveland Fed nowcast (both
$p < 10^{-21}$).

\begin{table}[t]
\centering
\caption{Headline head-to-head, all 29 paired events
(target months Jan 2025 -- Apr 2026, mixed MoM + YoY).  Errors are
in percentage points of CPI.  Paired $t$-test on absolute errors
($H_0$: equal MAE).}
\label{tab:headline}
\begin{tabular}{lrrrr}
\toprule
                   & $N$ & MAE (pp) & RMSE (pp) & bias (pp) \\
\midrule
Polymarket         & 29 & 0.124 & 0.225 & $+0.018$ \\
Cleveland Fed      & 29 & 0.124 & 0.174 & $-0.056$ \\
\midrule
Paired $t$-test    & \multicolumn{4}{c}{$t = 0.02$, $p = 0.99$} \\
Pearson $r$ vs.\ realised & \multicolumn{4}{c}{Poly $0.984$ \quad CF $0.992$} \\
\bottomrule
\end{tabular}
\end{table}

\subsection{Split by forecast horizon}

The headline tie hides a clear cross-pattern by horizon
(Table~\ref{tab:by_kind}).  On 13 monthly (MoM) prints Polymarket
beats the model by 4\,bp of MAE; on 16 annual (YoY) prints the model
beats Polymarket by 4\,bp of MAE.  Both differences are within the
noise floor of our sample.

\begin{table}[t]
\centering
\caption{Split by forecast kind.  Polymarket has the edge on monthly
(MoM) prints; Cleveland Fed has the edge on annual (YoY) prints.
Neither cross is statistically significant at $N$ this small, but the
sign is consistent.}
\label{tab:by_kind}
\begin{tabular}{llrrr}
\toprule
Kind & Forecaster & $N$ & MAE (pp) & RMSE (pp) \\
\midrule
\multirow{2}{*}{Monthly (MoM)} & Polymarket    & 13 & 0.080 & 0.106 \\
                               & Cleveland Fed & 13 & 0.123 & 0.171 \\
\midrule
\multirow{2}{*}{Annual (YoY)}  & Polymarket    & 16 & 0.160 & 0.288 \\
                               & Cleveland Fed & 16 & 0.124 & 0.177 \\
\bottomrule
\end{tabular}
\end{table}

\subsection{The truncated-bracket artefact}

Polymarket's RMSE is materially worse than its MAE
($0.225$ vs.\ $0.124$\,pp).  The single biggest contributor is the
March 2026 YoY event (id 208170), whose bracket grid was a set of
seven point brackets at $\{2.1\%, 2.2\%, 2.3\%, 2.4\%, 2.5\%, 2.6\%,
2.7\%\}$ -- so the implied mean was structurally bounded by the
$[2.1, 2.7]$\,pp interval -- while the realised CPI YoY print came
in at $3.26\%$, well outside the grid.  The implied mean was $2.32\%$,
off by $-0.93$\,pp.  The Polymarket operator subsequently added a
``Higher Brackets'' event (id 261155) covering
$\{2.7\%, 2.8\%, 2.9\%, 3.0\%, 3.1\%, 3.2\%, 3.3\%\}$; that event's
T-1 implied mean was $3.258\%$ -- within $0.002$\,pp of the realised
value.

The pattern is more general: the implied mean of a bracketed market
is a $\sigma$-truncated estimator and is biased toward the bracket
midrange whenever the realised value falls outside the bracket grid.
Operators can (and on Polymarket do) mitigate this by re-listing with
a wider grid, but doing so creates two competing events for the same
target, and the user faces a choice about which one to ``believe''.
We report both ``naive headline'' and ``addendum-preferring dedupe''
metrics; see Table~\ref{tab:robustness}.

\begin{table}[t]
\centering
\caption{Robustness checks.  ``Drop-1'' removes the single event with
the largest $\big||\varepsilon_{\text{poly}}| - |\varepsilon_{\text{cf}}|\big|$
(in both cases this is the March 2026 truncated-bracket event).
``Dedupe'' keeps the operator-issued addendum event when one exists.}
\label{tab:robustness}
\begin{tabular}{lrrrrr}
\toprule
                       & $N$ & MAE$_{\text{poly}}$ & MAE$_{\text{cf}}$ & RMSE$_{\text{poly}}$ & RMSE$_{\text{cf}}$ \\
\midrule
Headline (all events)  & 29 & 0.124 & 0.124 & 0.225 & 0.174 \\
Drop-1 (largest diff)  & 28 & 0.095 & 0.114 & 0.147 & 0.161 \\
Dedupe (prefer addendum) & 27 & \textbf{0.078} & 0.117 & \textbf{0.103} & 0.164 \\
\bottomrule
\end{tabular}
\end{table}

On the dedupe sample, the paired $t$-test on absolute errors gives
$t = -1.72, p = 0.097$ in favour of Polymarket -- a marginal effect
that should not be over-interpreted at $N = 27$, but is the opposite
direction of the headline tie.

\subsection{Volume tier}

Splitting the sample at the median Polymarket dollar volume per event
(\$591k), Polymarket's MAE is essentially the same in the two tiers
($0.119$\,pp above the median, $0.130$\,pp below).  The Cleveland Fed
nowcast model has \emph{lower} MAE on the low-volume tier
($0.114$\,pp), which makes sense: the low-volume tier is dominated by
the smaller, earlier 2025 markets when both forecasters had a stable
disinflation backdrop to lean on.  We do not find evidence that
``thicker markets are smarter markets'' in this sample.

\subsection{Calibration of probabilistic forecasts}

Polymarket emits a full probability distribution over brackets, not
just a point forecast.  The mean multi-class Brier score on 26 events
is $0.539$, compared to $0.80$ for a uniform-over-5-brackets
predictor.  This says Polymarket's distribution is informative, but
not sharply so -- much of the score comes from putting non-trivial
probability mass on brackets adjacent to the realised one, which is
weighed equally with putting mass on far brackets.  The worst
per-event Brier score is $1.79$ for November 2025 MoM, where realised
inflation came in at $0.02\%$ against an implied mean of $0.29\%$
(a surprise on which the Cleveland Fed model also failed, missing by
$+0.30$\,pp).

We also have six early-2024 legacy binary ``Will CPI MoM be $> X$\%?''
markets that did not survive into the bracketed era; their mean
binary Brier score is $0.217$ (against $0.25$ for a coin flip on
balanced YES/NO outcomes).  We exclude them from the main tables
because they cannot be compared apples-to-apples with the bracket
distributions of the modern markets.

\begin{figure}[t]
\centering
\includegraphics[width=\linewidth]{figure_1.png}
\caption{Forecast versus realised CPI for the 29 paired events.
\textbf{(a)} Polymarket bracket-implied mean; \textbf{(b)} Cleveland
Fed daily nowcast on the same cutoff date.  Annual (YoY) prints in
salmon circles, monthly (MoM) prints in blue squares.  The
$45^\circ$ line is the perfect-predictor reference.  Both forecasters
cluster tightly around the line; Polymarket's single visible outlier
in panel (a) at $(2.32, 3.26)$ is the truncated-bracket March 2026
event discussed in Section 5.3.}
\label{fig:scatter}
\end{figure}

\section{Discussion}

The most accurate one-line summary is: \emph{at the T-1 day horizon
for US CPI prints between January 2025 and April 2026, the median
Polymarket bracket-implied forecast and the Federal Reserve Bank of
Cleveland's daily nowcast were essentially tied in accuracy, both
with mean absolute error of about $0.12$\,pp.}  That is mildly
surprising on the prior that nowcasting models exploit a richer
information set -- BLS PPI components, energy spot prices, scanner
data -- than retail-and-prop traders are likely to.

A few sub-results push gently against the null:

\textbf{Polymarket wins on monthly MoM, loses on annual YoY.}
We did not predict this.  One speculative reading is that
annual YoY inflation is mechanically dominated by base effects (the
12-month-ago print) that a model can compute exactly and an
average trader has to look up; whereas monthly MoM noise depends more
on idiosyncratic shocks (rent surveys, holiday seasonality, energy
spot prices) where the marginal Polymarket trader may have an edge.
This would be worth testing with more data.

\textbf{Bracket design matters.}  The single largest Polymarket error
in our sample is a market-design failure, not a trader failure.  When
the operator listed a wider bracket grid for the same target,
Polymarket nailed it to within $0.002$\,pp.  This suggests two
operator-side levers: (i) wider default bracket grids, accepting some
loss of single-bracket pricing resolution; and (ii) automatic
addendum listing when the realised distribution starts to clip a
boundary bracket.

\textbf{The point-vs-distribution gap.}  Polymarket emits a full
distribution over brackets; we collapse it to a mean for direct
comparison with the Cleveland Fed nowcast (which is a point
estimate).  A richer benchmark would include the Cleveland Fed's
nowcast \emph{uncertainty}, which is published, or a distributional
nowcast such as those in \citet{knotek2017nowcasting}.  Pushing
distributional metrics like CRPS would also be informative.

\textbf{Calibration is OK, sharpness is not.}  A mean Brier of
$0.539$ on 5-bracket events is better than uniform ($0.80$) but
well above what a near-perfect probabilistic forecaster would
achieve.  This is broadly consistent with \citet{page2013calibration},
who find that prediction markets are reasonably well calibrated at
short horizons but not sharply so.

\section{Limitations}

\paragraph{Sample size.}  Twenty-nine paired observations is small.
The headline $p$-value of $0.99$ for ``no difference in MAE'' is
genuinely consistent with no difference, but it is also consistent
with effect sizes of $\pm 0.05$\,pp in either direction -- comparable
to the cross-pattern by horizon we observe.

\paragraph{Polymarket's CLOB look-back window.}  The
\texttt{/prices-history} endpoint rejects \texttt{(startTs, endTs)}
windows wider than $\sim$15 days at 1-hour fidelity.  We worked around
this by anchoring queries to the 14 days immediately before each
resolution; this is adequate for a T-1 day analysis but precludes
e.g.\ a T-30 or T-90 day forecast horizon study, which would more
fairly test the markets' aggregation properties.

\paragraph{Choice of economist benchmark.}  We use the Cleveland Fed
nowcasting model because it publishes daily and historical data is
recoverable as JSON.  The literal interpretation of ``economists''
might prefer the Bloomberg or Wall Street Journal CPI consensus, but
those are paywalled; the Philadelphia Fed Survey of Professional
Forecasters publishes quarterly and does not cover individual monthly
prints \citep{croushore2001realtime}.  The Cleveland Fed model has
also been shown to outperform the SPF consensus
\citep{knotek2017nowcasting}, so beating it is at least a meaningful
standard.

\paragraph{Bracket midpoint imputation.}  Treating each bracket as if
it had a point mass at the bracket midpoint biases the implied mean
when the within-bracket distribution is skewed, and is meaningless for
the open-ended end brackets ($\le v$ and $\ge v$) other than by
convention.  We assigned $\pm 0.05$\,pp offsets to those brackets,
matching the within-bracket grid spacing; sensitivity to this choice
is small for samples in which the boundary brackets carry little
probability and large when they do not, the latter being precisely the
truncation case discussed in Section 5.3.

\paragraph{No survival adjustment.}  We compare only events that
resolved.  Polymarket has not delisted any of our target inflation
markets, so there is no obvious survival bias in our sample, but we
cannot rule it out for the broader population of inflation-related
markets.

\paragraph{Single forecaster horizon.}  We compute everything at
T-1 day.  The accuracy of either forecaster at longer horizons (a
week, a month, a year) might tell a different story.

\section{Conclusion}

On 29 paired US CPI events from Polymarket between January 2025 and
April 2026, the bracket-implied point forecast tied the Cleveland Fed
nowcasting model on mean absolute error
($0.124$\,pp each, paired $t$-test $p = 0.99$).  The headline tie is
fragile: a single market-design artefact (a truncated bracket grid
on March 2026 YoY inflation, off by $-0.93$\,pp) accounts for almost
all of Polymarket's RMSE excess.  On a sample that prefers the
operator-issued ``Higher Brackets'' addendum when one exists,
Polymarket has the lower MAE ($0.078$ vs $0.117$\,pp).
Polymarket beats the model on monthly MoM prints and loses on annual
YoY prints; we do not have a confident story about why.

We make no triumphant claim.  At this sample size, the most defensible
statement is that liquid Polymarket inflation markets are roughly
competitive with a publicly-documented professional nowcasting system
-- which is itself a strong baseline against survey consensus
\citep{knotek2017nowcasting}.  This is consistent with the older
prediction-market literature \citep{wolfers2004prediction,
snowberg2013prediction} but extends it to a domain (high-frequency
US macro releases) where individual traders might plausibly have been
overmatched by the institutional forecasters.

\paragraph{Reproducibility.}  All numerical results in this paper are
generated by the scripts in \texttt{scripts/}; running
\texttt{python3 scripts/fetch\_data.py} (one-time, $\sim$3 min over
Polymarket's public APIs) followed by
\texttt{python3 scripts/analysis.py} and
\texttt{python3 scripts/figure\_1.py} reproduces every number and
the figure.

\bibliographystyle{plainnat}
\bibliography{references}

\end{document}
